\chapter{Reference: \code{print} (one-shot)}
\label{ch:refman-print}

\begin{aside}
\code{maya::print} renders a single element to stdout
inline. No alt-screen, no raw mode, no event loop. The
output joins the terminal's normal scrollback and the
function returns.
\end{aside}

\section{Signature}

\begin{lstlisting}[language=C++]
namespace maya {
    void print(Element e);
}
\end{lstlisting}

\section{Working example}

\begin{lstlisting}[language=C++]
#include <maya/maya.hpp>

using namespace maya;
using namespace maya::dsl;

int main() {
    print(
        v(
            t<"build complete"> | Bold | Fg<100, 220, 120>,
            blank_,
            text("warnings: 0") | Dim,
            text("errors:   0") | Dim
        ) | pad<1> | border_<Round>
    );

    print(t<"goodbye"> | Dim);
    return 0;
}
\end{lstlisting}

\section{When to use it}

\begin{itemize}[leftmargin=*]
  \item Build tool summaries.
  \item CLI status reports.
  \item One-off cards in scripts.
  \item Anywhere you'd otherwise write a sequence of
    \code{std::cout} lines.
\end{itemize}

\section{Notes}

\begin{itemize}[leftmargin=*]
  \item Each call writes a complete frame and a newline; the
    next call starts fresh below.
  \item Style escapes are emitted only if the terminal
    supports them; piping to a file strips them.
  \item No state survives between calls.
\end{itemize}

\section{Recap}

\begin{recap}
\begin{itemize}[leftmargin=*]
  \item One element in, one frame out.
  \item Inline; no alt-screen.
  \item Perfect for CLI output.
\end{itemize}
\end{recap}

\section*{Workshop}

\begin{enumerate}[leftmargin=*]
  \item Replace a multi-line \code{printf} with \code{print}
    in an existing tool.
  \item Build a "summary card" helper that takes a title and
    a list of pairs.
  \item Pipe the output to \code{cat} and confirm the
    escapes are gone.
\end{enumerate}
